\section{Survey Materials}
\label{sec:supp_materials}

This section reproduces the verbatim items administered to participants. Items are organized by the cluster organization used in the main text. All attitudinal items used a 7-point Likert scale anchored at 1 = \textit{Strongly disagree} and 7 = \textit{Strongly agree}, except where noted. Items marked (R) were reverse-scored before composite construction.

\subsection{Punitiveness Items}

The punitiveness aggregate composite combines three multi-item attitudinal subcomposites (Punish More, Parsimony, Three Strikes), two single-item attitudinal items (LWOP, Death Penalty), and the within-vignette $z$-scored Sentencing Decision. The 8-item attitudinal-only composite drops the sentencing decision and is reported in the main text wherever that decision is itself the dependent variable.

\textit{Punish More.} \textbf{(1)} The current level of punishment in the United States is too high. (R) \textbf{(2)} The solution to crime is more punishment.

\textit{Parsimony.} \textbf{(1)} Criminal punishments should be as light as possible. (R) \textbf{(2)} The state should not punish more than is absolutely necessary. (R)

\textit{Three Strikes.} \textbf{(1)} People convicted three times for a felony should be imprisoned for life. \textbf{(2)} People convicted three times for shoplifting should serve at least 25 years in prison.

\textit{LWOP.} Convicted murderers should never be allowed out of prison.

\textit{Death Penalty.} How much do you support or oppose the use of the death penalty to punish people convicted of first-degree murder? (5-point scale: 1 = \textit{Strongly oppose} to 5 = \textit{Strongly support}; rescaled to the 1--7 range for composite construction.)

\subsection{Crime Concerns Cluster}

\textit{Perceived Crime Rates.} \textbf{(1)} How serious is the problem of crime in the United States? (1 = \textit{Not at all serious}; 7 = \textit{Extremely serious}.) \textbf{(2)} Do you feel that violent crime just keeps rising?

\textit{Fear of Crime.} \textbf{(1)} How safe do you feel walking at night in the vicinity of your home? (R) \textbf{(2)} How much do you fear that you will be a victim of a serious crime? \textbf{(3)} How much do you fear that someone will break into your home in the next 12 months?

\subsection{Emotions Cluster}

\textit{Hatred.} \textbf{(1)} I feel hatred towards criminals. \textbf{(2)} I am sympathetic towards criminals. (R) \textbf{(3)} I feel that criminals are detestable human beings.

\textit{Anger.} \textbf{(1)} I feel anger towards criminals. \textbf{(2)} The crime problem makes me feel angry.

\subsection{Hostile Aggression Cluster}

\textit{Social Exclusion.} \textbf{(1)} By breaking society's laws, criminals exclude themselves from membership in society. \textbf{(2)} Having chosen to break the law, criminals do not deserve to enjoy the benefits that society provides its members. \textbf{(3)} Punishment communicates to criminals that they have no place in society.

\textit{Degradation.} \textbf{(1)} Regardless of their wrongdoing, prisoners deserve to be treated with dignity and respect. (R) \textbf{(2)} Punishment communicates to criminals that society has low regard for them. \textbf{(3)} Punishment should humiliate the criminal.

\textit{Infliction of Suffering.} \textbf{(1)} Infliction of suffering is a permissible goal of criminal punishment. \textbf{(2)} Punishment without suffering is not punishment.

\textit{Prison Violence Tolerance.} \textbf{(1)} I feel bad when a prisoner is assaulted in prison. (R) \textbf{(2)} Sexual assault of prisoners is a serious societal problem. (R)

\textit{Harsh Prison Conditions.} \textbf{(1)} Prison conditions should be extremely harsh. \textbf{(2)} Prisoners deserve decent living conditions. (R) \textbf{(3)} Our prisons are too comfortable.

\textit{Revenge.} \textbf{(1)} Society has the right to take revenge on criminal offenders. \textbf{(2)} Society should punish to get back at criminal offenders. \textbf{(3)} A father should not be prosecuted for killing the man who raped his 9-year-old daughter.

\subsection{Personality and Ideology Cluster}

\textit{Right-Wing Authoritarianism (RWA).} \textbf{(1)} What our country needs most is discipline, with everyone following our leaders in unity. \textbf{(2)} God's laws about abortion, pornography, and marriage must be strictly followed before it is too late. \textbf{(3)} There is nothing wrong with premarital sexual intercourse. (R) \textbf{(4)} Our society needs tougher government and stricter laws. \textbf{(5)} To preserve law and order, we should crack down harder on troublemakers.

\textit{Social Dominance Orientation (SDO).} \textbf{(1)} An ideal society requires some groups to be on top and others to be on the bottom. \textbf{(2)} Some groups of people are simply inferior to other groups. \textbf{(3)} No one group should dominate in society. (R) \textbf{(4)} Groups at the bottom of the social order are just as deserving as groups at the top. (R) \textbf{(5)} Group equality should not be our primary goal. \textbf{(6)} It is unjust to try to make groups equal. \textbf{(7)} We should do what we can to equalize conditions for different social groups. (R) \textbf{(8)} We should work to give all groups an equal chance to succeed. (R)

\textit{Trait Vengefulness.} \textbf{(1)} It is not worth my time or effort to pay back someone who has wronged me. (R) \textbf{(2)} It is important for me to get back at people who have hurt me. \textbf{(3)} There is nothing wrong in getting back at someone who has hurt you. \textbf{(4)} I don't just get mad, I get even. \textbf{(5)} I am not a vengeful person. (R)

\textit{Violence Proneness.} \textbf{(1)} A person has the right to kill to defend their property. \textbf{(2)} Physical punishment is necessary for raising children. \textbf{(3)} There is nothing wrong in torturing terrorists to obtain information about the threat they may pose to society. \textbf{(4)} Imagine that the United States is at war with a foreign country. After months of heavy losses, the US ground invasion has ground to a halt, with experts estimating an additional 20{,}000 US military fatalities. The US could end the conflict immediately by deploying a single nuclear weapon on the enemy capital, killing roughly 200{,}000 civilians. Should the US deploy the nuclear weapon?

\textit{Racial Resentment.} \textbf{(1)} Generations of slavery and discrimination have created conditions that make it difficult for Blacks to work their way out of the lower class. (R) \textbf{(2)} It's really a matter of some people not trying hard enough; if Blacks would only try harder they could be as well off as Whites. \textbf{(3)} Irish, Italian, Jewish, and many other minorities overcame prejudice and worked their way up. Blacks should do the same without any special favors. \textbf{(4)} Over the past few years, Blacks have gotten less than they deserve. (R)

\textit{Blood Sports Viewership.} Stem: ``How much do you like watching the following types of shows on TV?'' (1 = \textit{Hate it} to 7 = \textit{Love it}.) Items: Comedy; Swimming; Basketball; Boxing; Soccer (MLS, NWSL); Mixed Martial Arts (UFC); News; Hunting; Professional Wrestling (WWE, WWF). Composite combines the four blood-sports items: Boxing, MMA/UFC, Hunting, and Professional Wrestling. Other items were retained as filler/exploratory.

\subsection{Process Violations Cluster (Exploratory)}

This cluster was retained as exploratory after pre-registration. It was not part of the confirmatory tests of H2 but is included in correlational analyses for completeness.

\textit{Due Process.} \textbf{(1)} The legal system's loopholes and technicalities make it too difficult to convict criminals for their crimes. \textbf{(2)} No matter how abhorrent the crime is, every defendant deserves the protection of due process of law. (R) \textbf{(3)} People charged with committing violent crimes should be convicted even if police investigators cut some corners.

\textit{Uncertain Evidence.} \textbf{(1)} Jailhouse informants can be trusted when they testify that a cell-mate had confessed to committing a crime. \textbf{(2)} Generally, the police make an arrest only when they are sure who committed the crime. \textbf{(3)} When it is the suspect's word against the police officer's word, we should believe the police. \textbf{(4)} If police officers know that a defendant is guilty, it is okay for them to lie when testifying in court to help obtain a conviction.

\subsection{Sentencing Vignettes}
\label{sec:supp_vignettes}

Each participant was randomly assigned to one of three vignettes. All three resulted in a second-degree murder conviction. The framing instructions and dependent measure were identical across vignettes.

\textit{Framing instructions} (presented after each vignette):
\begin{quote}
You are now being asked what punishment Darryl deserves for this crime. Assume that Darryl's conduct amounts to second degree murder. The law states: ``Whoever is guilty of second degree murder shall be imprisoned for any term of years, or for life in prison.'' This means that you can decide on a sentence ranging from zero years in prison up to 50 years in prison (which is generally equivalent to life in prison). Most sentences given by judges for similar charges range from 20 to 30 years.
\end{quote}

\textit{Dependent measure}: ``What sentence would you give Darryl Smith for killing Frankie Williams? I would sentence Darryl to this number of years in prison: (for a life sentence, mark 50 years).''

\textit{Vignette 1: Stranger Felony-Murder} ($n = 168$). Darryl Smith had just lost his job for the second time due to closures of nearby factories. Days later, around 9:30~pm, Darryl followed Frankie Williams, a teacher, as she was walking towards the parking lot of a shopping mall. Darryl snuck up on Frankie as she turned a dark corner and tried to snatch her purse. Frankie managed to hold onto her purse and swung her heavy shopping bag at him, causing a deep cut above his eyebrow. With blood pouring into his eye, Darryl became furious. He twisted Frankie's arm, snatched the purse, and pushed her against a railing. Frankie fell over the railing, struck her head on the pavement below, and died of her injuries.

\textit{Vignette 2: Domestic Violence Murder} ($n = 176$). Darryl Smith and Frankie Williams had been dating on and off for some 2 years. Throughout, Darryl's behavior was volatile and occasionally also violent. Frankie had twice filed reports of domestic abuse with the police and had recently obtained a restraining order against Darryl. Just weeks after he was presented with the restraining order, Darryl burst into Frankie's apartment. His breath smelled of alcohol. He stated that he wanted just to talk to straighten things out. Frankie demanded that he leave. When he refused, she ran toward the door. Darryl gave chase. As Frankie ran into the street, she was struck by oncoming traffic and died of her injuries.

\textit{Vignette 3: Organized Crime Murder} ($n = 152$). Darryl Smith ran an illegal operation of stealing luxury cars and selling them for parts. At some point, Darryl started suspecting that Frankie Williams, a local drug dealer, was snitching on him to the police. Darryl confronted Frankie about his connections to the police. Frankie responded aggressively that he'd make sure that Darryl ends up in prison. Soon thereafter, Frankie was spotted near Darryl's garage taking photographs of the stolen cars. Darryl decided to beat Frankie up in order to teach him a lesson. Together with an accomplice, Darryl ambushed Frankie and beat him with a baseball bat. Frankie died from his injuries the following day.

\subsection{Open-Ended Justifications}

After the sentencing decision, participants saw two open-ended prompts presented in sequence:

\begin{quote}
\textbf{Prompt 1.} Please explain why you recommended this sentence for Darryl.

\textbf{Prompt 2.} What do you hope this sentence will accomplish?
\end{quote}

The two responses were concatenated for NLP analysis (combined $M = 48.3$ words, $SD = 35.9$, $\mathit{Mdn} = 37$, range 5--231).

\subsection{Attention Checks and Demographics}

Two attention-check items were embedded in the survey: ``To make sure that the experiment is running correctly, please select the strongly agree option'' and ``What is 8 + 7 equal to?'' Failure on either resulted in exclusion (18 participants excluded; final $N = 496$).

Demographics: age (years, integer), gender, race/ethnicity, years of education, and political orientation (1 = \textit{Strongly liberal} to 7 = \textit{Strongly conservative}, recoded into a three-level political-group variable for moderation analyses).

\section{Pre-Registration Status and Deviations}
\label{sec:supp_prereg}

The study was pre-registered on AsPredicted (\#129382 supersedes the preliminary version; current version archived on the OSF). The pre-registration specified the three confirmatory hypotheses (H1, H2, H3), the punitiveness composite, the cluster organization, the vignette randomization, the within-vignette $z$-scoring of the sentencing decision, and a target sample size of 480 (achieved $N = 496$).

The following deviations from pre-registration are transparently reported here.

\textit{Cluster organization}. The pre-registration listed five clusters (Crime Concerns, Process Violations, Emotions, Hostile Aggression, Personality). The Process Violations cluster was demoted to exploratory status after pre-registration; the four remaining clusters constitute the confirmatory analysis. This decision was made before data analysis on the current sample. Process Violations is reported throughout the manuscript and in this appendix for completeness; it does not enter the central H2 test.

\textit{Item revisions}. The pre-registration listed three items per construct in places where the final survey used two (e.g., Suffering, Prison Violence). The reductions reflect substantive revisions during instrument finalization, not post-hoc pruning. All Cronbach's $\alpha$ estimates were computed on the items as administered (Table~\ref{tab:reliability}).

\textit{NLP analysis plan}. The pre-registration committed to a multi-method NLP approach including dictionary analysis, sentiment analysis, topic modeling, semantic similarity, zero-shot classification, and moral foundations analysis. The final pipeline implements all five core method classes. The Moral Foundations Dictionary was scored but did not produce findings beyond what the Empath \citep{fast2016empath} categories already capture; it is omitted from the main text and from this appendix to avoid redundancy. The pre-registered BART-MNLI classifier was retired after benchmark validation showed DeBERTa-v3-large \citep{he2023deberta} substantially outperformed it (87.5\% vs.\ 77.5\% top-1 accuracy on a 40-sentence ground-truth set; see Section~\ref{sec:supp_nlp_validation}). A forced-choice large-language-model classifier (Claude Haiku~4.5; \citealp{anthropicclaude2024}) was added as a third classification method.

\textit{Sample size}. The pre-registration set 480 as the target. The realized $N$ after attention-check exclusions was 496. No additional exclusions were applied.

\textit{Tautology sensitivity analyses, TOST equivalence tests, and the cultural-default multiverse} were not pre-registered but were added in response to peer-reader concerns. They are reported throughout this appendix as exploratory.
\section{Study 1 Supplementary Analyses}
\label{sec:supp_study1}

\subsection{Item-Level H1 Correlations With FDR}
\label{sec:supp_item_H1}

Hypothesis~1 was tested at the construct level in the main text (Table~\ref{tab:H1_correlations}). For completeness, Table~\ref{tab:item_h1} reports the same hypothesis at the item level, sorted by $|r|$. All 56 items show positive correlations with the punitiveness aggregate. Forty of 56 are significant after Benjamini--Hochberg FDR correction at $q = .05$.

\begin{table}[htbp]
  \centering
  \caption{Item-level Pearson correlations with the punitiveness aggregate, sorted by $|r|$. All correlations are positive. Item codes correspond to the export tags in the dataset and to the items listed in Section~\ref{sec:supp_materials}. $N = 496$.}
  \label{tab:item_h1}
  \footnotesize
  \begin{tabular}{lrrrl}
    \toprule
    Item & $r$ & $p$ & Sig & Construct \\
    \midrule
    exclusion\_2     & $.57$ & $<.001$ & *** & Social Exclusion \\
    rwa\_5           & $.56$ & $<.001$ & *** & RWA \\
    exclusion\_1     & $.51$ & $<.001$ & *** & Social Exclusion \\
    suffer\_1        & $.51$ & $<.001$ & *** & Infliction of Suffering \\
    harsh\_1         & $.50$ & $<.001$ & *** & Harsh Conditions \\
    revenge\_1       & $.49$ & $<.001$ & *** & Revenge \\
    hate\_1          & $.49$ & $<.001$ & *** & Hatred \\
    suffer\_2        & $.48$ & $<.001$ & *** & Infliction of Suffering \\
    rwa\_4           & $.48$ & $<.001$ & *** & RWA \\
    revenge\_2       & $.47$ & $<.001$ & *** & Revenge \\
    harsh\_3         & $.47$ & $<.001$ & *** & Harsh Conditions \\
    raceresent\_2    & $.45$ & $<.001$ & *** & Racial Resentment \\
    hate\_3          & $.45$ & $<.001$ & *** & Hatred \\
    \multicolumn{5}{c}{\textit{(43 additional items, all positive, range $|r| = .04$ to $.44$)}} \\
    venge\_3         & $.07$ & $.143$  & ns  & Trait Vengefulness \\
    venge\_4         & $.06$ & $.183$  & ns  & Trait Vengefulness \\
    venge\_1         & $.04$ & $.379$  & ns  & Trait Vengefulness (R) \\
    \bottomrule
  \end{tabular}

  \vspace{0.3em}
  \footnotesize
  \textit{Note.} Full 56-row table available in the project repository (\texttt{H1\_item\_level\_correlations.csv}). Forty items survive FDR correction at $q = .05$.
\end{table}

\subsection{Vignette Stability of Punitiveness Correlations}
\label{sec:supp_vignette_stab}

Each correlation reported in the main text is computed across all 496 participants. Because vignette assignment was randomized, between-vignette differences in the correlations should reflect sampling noise rather than systematic differences in the underlying psychology. Table~\ref{tab:vignette_stab} reports the punitiveness correlation for each construct, separately within each of the three vignettes plus the overall correlation. The pattern is highly stable: hostile aggression constructs (exclusion, suffering, harsh conditions, hatred) remain the strongest correlates within every vignette, and the rank ordering of constructs is preserved.

\begin{table}[htbp]
  \centering
  \caption{Pearson correlation between each correlate construct and the punitiveness aggregate, within each vignette and overall. The Range column is the maximum minus the minimum of the three vignette-specific correlations. Constructs are sorted by overall $r$. $n = 168$ for Vignette~1, $n = 176$ for Vignette~2, $n = 152$ for Vignette~3, $N = 496$ overall.}
  \label{tab:vignette_stab}
  \footnotesize
  \begin{tabular}{lcccc c}
    \toprule
    Construct & Vignette 1 & Vignette 2 & Vignette 3 & Overall & Range \\
    \midrule
    Harsh Conditions   & $.62$ & $.54$ & $.49$ & $.56$ & $.13$ \\
    Exclusion          & $.64$ & $.46$ & $.51$ & $.54$ & $.18$ \\
    Hatred             & $.60$ & $.49$ & $.49$ & $.53$ & $.11$ \\
    Suffering          & $.59$ & $.51$ & $.43$ & $.52$ & $.16$ \\
    RWA                & $.60$ & $.43$ & $.46$ & $.50$ & $.17$ \\
    Racial Resentment  & $.51$ & $.40$ & $.49$ & $.46$ & $.11$ \\
    Anger              & $.55$ & $.39$ & $.38$ & $.45$ & $.17$ \\
    Degradation        & $.52$ & $.41$ & $.38$ & $.44$ & $.14$ \\
    Revenge            & $.51$ & $.40$ & $.37$ & $.43$ & $.14$ \\
    Due Process        & $.45$ & $.45$ & $.39$ & $.43$ & $.06$ \\
    Violence Proneness & $.47$ & $.33$ & $.43$ & $.41$ & $.14$ \\
    Crime Rates        & $.41$ & $.42$ & $.36$ & $.40$ & $.06$ \\
    Uncertain Evidence & $.29$ & $.37$ & $.42$ & $.35$ & $.13$ \\
    Prison Violence    & $.31$ & $.39$ & $.35$ & $.35$ & $.08$ \\
    SDO                & $.36$ & $.32$ & $.28$ & $.32$ & $.08$ \\
    Fear               & $.20$ & $.30$ & $.03$ & $.18$ & $.27$ \\
    Vengefulness       & $.21$ & $.05$ & $.09$ & $.12$ & $.16$ \\
    Blood Sports       & $.12$ & $.10$ & $.12$ & $.11$ & $.02$ \\
    \bottomrule
  \end{tabular}
\end{table}

\subsection{Sentence Anchor Analysis}
\label{sec:supp_anchor}

The vignette framing instructions described typical sentences as ranging from 20 to 30 years, anchoring the midpoint at 25 years. Whether participants converged on the midpoint or systematically departed from it bears on whether the sentencing decision is informative beyond the framing.

Single-sample $t$-tests against the 25-year midpoint show that participants pooled across vignettes recommended substantially longer sentences than the anchor midpoint, $M = 35.19$ years, $t(495) = 18.14$, $p < .001$, Cohen's $d = 0.81$. The departure replicates within each vignette (Table~\ref{tab:anchor}). The Wilcoxon signed-rank test, robust to non-normality, confirms the same pattern at $p < .001$ in all four tests. A signed-distance analysis shows that the magnitude of departure from the midpoint correlates with hostile aggression ($r = .26$, $p < .001$) and punitiveness ($r = .49$, $p < .001$): participants high in punitive dispositions departed from the anchor in the punitive direction more than the average departure already shown by the sample as a whole.

\begin{table}[htbp]
  \centering
  \caption{Single-sample $t$-tests of the sentencing decision against the 25-year anchor midpoint, separately by vignette and pooled. $d$ is Cohen's $d$. Wilcoxon $W$ is the signed-rank test statistic; both Wilcoxon $p$-values were $< .001$ for all rows.}
  \label{tab:anchor}
  \small
  \begin{tabular}{lrrrrrr}
    \toprule
    Vignette & $n$ & $M$ & $SD$ & $M - 25$ & $t$ & $d$ \\
    \midrule
    Vignette 1 & 168 & $34.41$ & $12.59$ & $9.41$  & $9.69$  & $0.75$ \\
    Vignette 2 & 176 & $37.46$ & $12.60$ & $12.46$ & $13.12$ & $0.99$ \\
    Vignette 3 & 152 & $33.41$ & $11.99$ & $8.41$  & $8.64$  & $0.70$ \\
    Pooled     & 496 & $35.19$ & $12.51$ & $10.19$ & $18.14$ & $0.81$ \\
    \bottomrule
  \end{tabular}
\end{table}

\subsection{Bootstrap CIs for the Hostile Aggression Advantage}
\label{sec:supp_bootstrap}

Figure~\ref{fig:bootstrap} in the main text presents the BCa 95\% confidence intervals for the difference between $r_\textit{hostile}$ and $r_\textit{crime concerns}$ at three stringency levels of construct inclusion. Table~\ref{tab:supp_bootstrap_full} reports the same intervals plus the parallel intervals for the three other clusters tested in the H2 extension: emotions, personality, and process violations. All five rows show CIs that exclude zero.

\begin{table}[htbp]
  \centering
  \caption{Bootstrap 95\% confidence intervals (BCa, 10{,}000 iterations) for the difference between the cluster correlation with punitiveness and the crime concerns correlation with punitiveness ($r = .322$). Dependent-correlations bootstrap was conducted on a single complete-cases dataset for direct comparability. $N = 496$ for all rows.}
  \label{tab:supp_bootstrap_full}
  \small
  \begin{tabular}{lrrrr}
    \toprule
    Comparison & Observed $\Delta r$ & Boot $M$ & BCa Lower & BCa Upper \\
    \midrule
    Hostile Aggression --- 6 constructs (Original)   & $.264$ & $.264$ & $.174$ & $.367$ \\
    Hostile Aggression --- 4 constructs (Conservative) & $.226$ & $.225$ & $.134$ & $.331$ \\
    Hostile Aggression --- 2 constructs (Strictest)   & $.221$ & $.222$ & $.125$ & $.323$ \\
    Emotions cluster                                 & $.213$ & $.213$ & $.117$ & $.320$ \\
    Personality cluster                              & $.153$ & $.153$ & $.054$ & $.261$ \\
    \bottomrule
  \end{tabular}
\end{table}

\subsection{Parsimony Sensitivity}
\label{sec:supp_parsimony}

The parsimony subscale ($\alpha = .48$) was retained in the punitiveness composite per pre-registration but, given its low reliability, the central H2 test was re-run with parsimony excluded. Removing parsimony reduces the punitiveness composite to a 6-item attitudinal-only construct. The hostile aggression advantage is preserved: $r_\textit{hostile} = .58$, $r_\textit{crime} = .33$, $\Delta r = .25$, Steiger's $Z = 5.87$, $p < .001$. The original $\Delta r = .264$ and the parsimony-removed $\Delta r = .247$ differ by less than $0.02$, well within the bootstrap CI of either. Parsimony's low reliability is therefore not driving the central finding.

\subsection{Confirmatory Factor Analysis}
\label{sec:supp_cfa}

A four-factor CFA model loaded the 16 confirmatory cluster constituents on Crime Concerns, Emotions, Hostile Aggression, and Personality \& Ideology factors. All 16 standardized loadings were positive and significant, $\lambda$s ranging from $.37$ for blood sports viewership to $.91$ for hatred (Table~\ref{tab:cfa_loadings}). Inter-factor correlations were uniformly large, with the strongest links between Emotions and Hostile Aggression ($r = .80$) and between Hostile Aggression and Personality \& Ideology ($r = .78$; Table~\ref{tab:cfa_corrs}).

Global model fit fell short of conventional thresholds: $\chi^2(98) = 699.10$, $p < .001$, CFI $= .858$, TLI $= .826$, RMSEA $= .111$ [.104, .119], SRMR $= .069$. The imperfect fit indicates that the four-cluster organization is best understood as a conceptual scaffold rather than a strict latent measurement model, as discussed in the main text. The substantive findings, which operate at the level of observed cluster composites, are not affected by this measurement-model limitation.

A five-factor model that separated Process Violations from Personality \& Ideology was estimated for comparison (Table~\ref{tab:cfa_compare}). The five-factor model showed slightly better RMSEA but identical CFI/TLI to the four-factor model, with substantially worse AIC and BIC. The two solutions are statistically close; the four-factor structure is preferred for parsimony and theoretical clarity.

\begin{table}[htbp]
  \centering
  \caption{Standardized loadings from the four-factor CFA. All $p < .001$.}
  \label{tab:cfa_loadings}
  \small
  \begin{tabular}{lll r r}
    \toprule
    Factor & Indicator & Loading & $SE$ & $p$ \\
    \midrule
    Crime Concerns       & Crime Rates       & $.81$ & $.06$ & $<.001$ \\
    Crime Concerns       & Fear              & $.48$ & $.05$ & $<.001$ \\
    Emotions             & Hatred            & $.91$ & $.02$ & $<.001$ \\
    Emotions             & Anger             & $.79$ & $.02$ & $<.001$ \\
    Hostile Aggression   & Exclusion         & $.81$ & $.02$ & $<.001$ \\
    Hostile Aggression   & Degradation       & $.83$ & $.02$ & $<.001$ \\
    Hostile Aggression   & Suffering         & $.83$ & $.02$ & $<.001$ \\
    Hostile Aggression   & Prison Violence   & $.55$ & $.03$ & $<.001$ \\
    Hostile Aggression   & Harsh Conditions  & $.84$ & $.02$ & $<.001$ \\
    Hostile Aggression   & Revenge           & $.75$ & $.02$ & $<.001$ \\
    Personality          & RWA               & $.71$ & $.03$ & $<.001$ \\
    Personality          & SDO               & $.68$ & $.03$ & $<.001$ \\
    Personality          & Vengefulness      & $.37$ & $.04$ & $<.001$ \\
    Personality          & Violence Proneness & $.78$ & $.02$ & $<.001$ \\
    Personality          & Racial Resentment & $.70$ & $.03$ & $<.001$ \\
    Personality          & Blood Sports      & $.38$ & $.04$ & $<.001$ \\
    \bottomrule
  \end{tabular}
\end{table}

\begin{table}[htbp]
  \centering
  \caption{Inter-factor correlations from the four-factor CFA. All $p < .001$.}
  \label{tab:cfa_corrs}
  \small
  \begin{tabular}{ll r}
    \toprule
    Factor 1 & Factor 2 & $r$ \\
    \midrule
    Crime Concerns & Emotions          & $.58$ \\
    Crime Concerns & Hostile Aggression & $.50$ \\
    Crime Concerns & Personality       & $.50$ \\
    Emotions       & Hostile Aggression & $.80$ \\
    Emotions       & Personality       & $.65$ \\
    Hostile Aggression & Personality   & $.78$ \\
    \bottomrule
  \end{tabular}
\end{table}

\begin{table}[htbp]
  \centering
  \caption{Comparison of four-factor and five-factor CFA solutions. Both models were fit with maximum likelihood estimation in \texttt{lavaan}.}
  \label{tab:cfa_compare}
  \small
  \begin{tabular}{lrrrrrr}
    \toprule
    Model & $\chi^2$ & $df$ & CFI & RMSEA & AIC & BIC \\
    \midrule
    4-factor & $699.10$ & $98$  & $.86$ & $.111$ & $25{,}210$ & $25{,}370$ \\
    5-factor & $805.81$ & $125$ & $.86$ & $.105$ & $27{,}819$ & $28{,}013$ \\
    \bottomrule
  \end{tabular}
\end{table}

\subsection{TOST Equivalence Testing}
\label{sec:supp_tost}

For the seven correlations that did not reach conventional significance in the main text, we conducted Two One-Sided Tests for equivalence \citep{lakens2017equivalence} with bounds $\pm 0.15$ (a small-effect benchmark). Five of the seven were statistically equivalent to zero by this test; two were not (Table~\ref{tab:tost}). For the two non-equivalent correlations, we re-ran the test at three alternative bounds ($\pm 0.10$, $\pm 0.125$, $\pm 0.20$) to characterize sensitivity. The voyage prosocial-dark gap correlation with hostile aggression and the dictionary prosocial correlation with hostile aggression both fail equivalence at $\pm 0.10$ and $\pm 0.15$ but are equivalent at $\pm 0.20$. We report these as null results consistent with the cultural-default reading, with the caveat that we cannot rule out small effects below $|r| = .15$.

\begin{table}[htbp]
  \centering
  \caption{TOST equivalence tests for non-significant correlations from the facade matrix. Equivalence bounds are $\pm 0.15$. ``Equivalent'' indicates $p_\textit{TOST} < .05$.}
  \label{tab:tost}
  \footnotesize
  \begin{tabular}{lrrlrrl}
    \toprule
    Correlation pair & $n$ & $r$ & 95\% CI & $p_\textit{NHST}$ & $p_\textit{TOST}$ & Equiv. \\
    \midrule
    BGE pro--dark gap $\times$ Hostile Aggr     & 496 & $-.04$ & $[-.13,\,.05]$ & $.41$ & $.006$ & Yes \\
    voyage pro--dark gap $\times$ Hostile Aggr  & 496 & $-.09$ & $[-.17,\,.00]$ & $.06$ & $.070$ & No \\
    mpnet pro--dark gap $\times$ Hostile Aggr   & 496 & $-.03$ & $[-.12,\,.05]$ & $.45$ & $.005$ & Yes \\
    VADER $\times$ Hostile Aggr                 & 496 & $-.05$ & $[-.14,\,.04]$ & $.25$ & $.014$ & Yes \\
    BGE pro--dark gap $\times$ Crime Concerns   & 496 & $.02$  & $[-.07,\,.11]$ & $.68$ & $.002$ & Yes \\
    voyage pro--dark gap $\times$ Crime Concerns & 496 & $.02$ & $[-.07,\,.11]$ & $.62$ & $.002$ & Yes \\
    Dictionary prosocial $\times$ Hostile Aggr  & 496 & $-.08$ & $[-.17,\,.01]$ & $.08$ & $.055$ & No \\
    \bottomrule
  \end{tabular}
\end{table}

\subsection{Demographic Moderation}
\label{sec:supp_demo_mod}

The demographic moderation analyses reported in the main text (Tables~\ref{tab:demographic_moderation_punit} and \ref{tab:demographic_moderation_sent}) tested whether age, gender, race, education, or political orientation moderated the association between each correlate cluster and punitiveness or the sentencing decision. Across 48 interaction tests, only one survived FDR correction: Crime Concerns $\times$ Age predicting punitiveness ($b = .16$, $p_\textit{FDR} = .000061$). Older participants showed a stronger crime-concerns--punitiveness link than younger participants. No other interaction reached even nominal significance after correction. The hostile aggression advantage holds across age, gender, race, education, and political orientation.
\section{Study 2 NLP Pipeline Details}
\label{sec:supp_nlp_pipeline}

\subsection{Text Preprocessing}

Both open-ended responses were retrieved from Qualtrics in raw form, then processed by the following pipeline. Step~1: Unicode normalization (NFKC) and replacement of non-ASCII whitespace with standard spaces. Step~2: lowercasing for downstream lexicon matching while retaining a cased copy for embedding models. Step~3: removal of HTML entities and surveyware artifacts. Step~4: concatenation of the two responses with a single space delimiter. Step~5: word count computation on the concatenated text, used as a control variable in regressions.

We did not stem, lemmatize, or remove stopwords. Stemming and lemmatization can change embeddings in unpredictable ways for short responses, and stopword removal would distort the dictionary count denominators. Length normalization was applied at the dictionary stage (counts divided by word count) and was not needed for the embedding methods.

\subsection{Dictionary Categories}
\label{sec:supp_dict_contents}

Ten justification categories were defined by manually curated keyword lists. Five were classified as conventionally prosocial (deterrence, incapacitation, rehabilitation, retribution, norm expression) and five as dark (revenge, suffering, degradation, exclusion, victim focus). Each list contained 5--13 stems and short multi-word phrases. Counts were computed as the number of word tokens whose start matched a single-token stem, plus the number of substring occurrences for multi-token phrases, divided by the response word count. Each response was assigned the category with the highest normalized count as its dictionary classification. Responses with zero matches in any category ($n = 75$) were left unclassified.

\begin{table}[htbp]
  \centering
  \caption{Dictionary contents for the ten justification categories.}
  \label{tab:supp_dict}
  \footnotesize
  \begin{tabular}{lp{0.72\textwidth}}
    \toprule
    Category & Keywords (single tokens are stem-matched; phrases in italics are substring-matched) \\
    \midrule
    \textit{Prosocial categories} \\
    Deterrence       & deter, prevent, discourage, warning, example, lesson, \textit{think twice}, \textit{stop others} \\
    Incapacitation   & protect, safety, safe, remove, dangerous, threat, public, \textit{keep away}, \textit{off the streets} \\
    Rehabilitation   & rehabilitate, reform, change, learn, treatment, therapy, help, improve, \textit{better person}, \textit{second chance}, program, education \\
    Retribution      & deserve, fair, just, justice, proportional, appropriate, consequences, accountable, responsible, \textit{pay for} \\
    Norm Expression  & wrong, unacceptable, society, values, moral, message, condemn, show, signal \\
    \addlinespace
    \textit{Dark categories} \\
    Revenge          & revenge, vengeance, payback, \textit{get back}, retaliate, \textit{eye for}, \textit{taste of} \\
    Suffering        & suffer, pain, hurt, agony, torment, misery, \textit{feel what}, rot \\
    Degradation      & humiliate, shame, disgrace, worthless, scum, animal, monster \\
    Exclusion        & remove, rid, eliminate, banish, \textit{no place}, \textit{throw away}, \textit{lock up forever} \\
    Victim Focus     & victim, family, closure, peace, \textit{loved ones} \\
    \bottomrule
  \end{tabular}
\end{table}

We did not develop or apply a separate moral-foundations dictionary in the main pipeline; preliminary scoring with the extended Moral Foundations Dictionary produced no findings beyond what Empath \citep{fast2016empath} already captured, so it was omitted to avoid redundant reporting.

\subsection{Zero-Shot Classifier (DeBERTa)}

The zero-shot classifier was DeBERTa-v3-large fine-tuned on natural-language-inference data \citep{he2023deberta}, accessed via the Hugging Face \texttt{zero-shot-classification} pipeline. The eight candidate labels were: \textit{deterrence and prevention}, \textit{incapacitation and public safety}, \textit{rehabilitation and reform}, \textit{proportional justice}, \textit{societal condemnation}, \textit{punishment and suffering}, \textit{revenge and payback}, and \textit{victim closure}. The pipeline returned a probability for each label and we used the highest-probability label as the response's primary classification. Multi-label scoring was used (probabilities do not sum to 1), so the absolute level of each label is interpretable as the model's confidence that the response falls in that category.

The eight-label set differs from the dictionary's ten categories in two ways. First, deterrence and incapacitation are merged into a single ``deterrence and prevention'' category, reflecting how those goals usually co-occur in everyday talk and how participants tended to invoke them in tandem. Second, ``norm expression'' is replaced by the more readable ``societal condemnation'', and ``victim focus'' by the more action-oriented ``victim closure''. The dictionary and embedding methods kept the original ten categories with deterrence and incapacitation distinct.

\subsection{Forced-Choice Classifier (Claude)}

The forced-choice classifier sent each response to Claude Haiku~4.5 \citep{anthropicclaude2024} with the same eight candidate labels. The prompt asked the model to return the single best-fitting label and a confidence score on a 1--5 scale, plus per-label scores on the same 1--5 scale. The full prompt is reproduced below; placeholders are shown in angle brackets.

\begin{quote}\small\itshape
You are a careful research assistant classifying short open-ended responses from a survey about criminal sentencing. The respondent has just chosen a sentence for a hypothetical defendant and has been asked to explain why.

Read the response below carefully and classify it according to the following eight justification categories: (1)~deterrence and prevention; (2)~incapacitation and public safety; (3)~rehabilitation and reform; (4)~proportional justice; (5)~societal condemnation; (6)~punishment and suffering; (7)~revenge and payback; (8)~victim closure.

Provide:
\begin{itemize}\itemsep -0.3em
  \item The single best-fitting category (one of the eight above).
  \item A confidence score on a 1--5 scale, where 5 = the response is unambiguously this category.
  \item A per-category score on a 1--5 scale for each of the eight categories, indicating how much the response expresses that category.
\end{itemize}

Return your answer as JSON in the following format: \texttt{\{\textit{top}: \textit{<category>}, \textit{confidence}: \textit{<1--5>}, \textit{scores}: \{\textit{<category>}: \textit{<1--5>}, ...\}\}}.

Response: $\langle$\textit{response\_text}$\rangle$
\end{quote}

We used the same eight-label set as DeBERTa for direct comparability. Per-category 1--5 scores were normalized to the 0--1 range for the multi-label vector that enters the prosocial--dark gap calculations.

\subsection{Embedding Models and Prototype-Similarity Method}
\label{sec:supp_embeddings}

Three sentence-embedding models were used to encode each response and each prototype: BGE-large-en-v1.5 \citep{xiao2024bge}, voyage-3-large \citep{voyageai2024}, and Sentence-MPNet-base-v2 \citep{song2020mpnet, reimers2019sentence}. The three models span architectural and training-data variation: BGE is a contrastively trained bidirectional encoder; voyage-3-large is a recent commercial embedding model; Sentence-MPNet is a 2021-era widely used baseline. Embeddings were L2-normalized before cosine similarity computation.

For each prototype theme, a centroid embedding was computed as the mean of its three prototype-sentence embeddings. Each response embedding was compared against the seven theme centroids by cosine similarity. The result is a 7-dimensional theme-similarity vector per response per embedding model.

Three aggregate prosocial-minus-dark gap measures were computed from each model:
\begin{itemize}\itemsep -0.3em
\item \textbf{B3a} (mean prosocial similarity minus mean dark similarity): \\
$\frac{1}{4}(\text{deterrence} + \text{incapacitation} + \text{rehabilitation} + \text{retribution}) - \frac{1}{3}(\text{revenge} + \text{suffering} + \text{exclusion})$.
\item \textbf{B3b} (max prosocial similarity minus max dark similarity).
\item \textbf{B3c} (one-hot top-similarity prosocial--dark indicator).
\end{itemize}
The main text uses B3a (the aggregate-mean version) throughout. B3b and B3c are reported in the project repository for sensitivity.

\subsection{Prototype Sentence Sets}
\label{sec:supp_prototypes}

Five prototype sentence sets were used, each containing seven categories with three sentences per category (105 prototype sentences total). The five sets vary in register and word choice while preserving the underlying conceptual targets. The Original set is from the project pipeline as initially developed; the Formal set uses academic/legal register; the Colloquial set uses lay register; the LayFormal set uses plain English while preserving the conceptual distinctness of incapacitation versus deterrence; the AntiTaut set was designed to minimize lexical overlap with the punitiveness composite items, addressing the concern that observed prosocial-dark gap correlations might reflect surface lexical similarity rather than thematic alignment.

The Original prototype sentences are reproduced below; the four alternative sets follow the same conceptual targets in different registers and are available in the project repository in full.

\begin{table}[htbp]
  \centering
  \caption{Original prototype sentences. Each theme has three prototypes; the centroid of the three serves as the theme's reference embedding.}
  \label{tab:supp_protos_orig}
  \footnotesize
  \begin{tabular}{lp{0.72\textwidth}}
    \toprule
    Theme & Prototype sentences \\
    \midrule
    Deterrence    & 1.~This sentence will deter others from committing similar crimes. \newline 2.~Punishment sends a message that crime does not pay. \newline 3.~This will make others think twice before breaking the law. \\
    \addlinespace
    Incapacitation & 1.~This keeps dangerous criminals away from society. \newline 2.~The public needs to be protected from this offender. \newline 3.~Society is safer with this person behind bars. \\
    \addlinespace
    Rehabilitation & 1.~Time in prison gives the offender a chance to reform. \newline 2.~With proper programs, this person can become a productive member of society. \newline 3.~Everyone deserves a second chance to become a better person. \\
    \addlinespace
    Retribution    & 1.~This sentence is proportional to the severity of the crime. \newline 2.~Justice requires that wrongdoers face appropriate consequences. \newline 3.~The offender deserves this punishment for their actions. \\
    \addlinespace
    Revenge        & 1.~He deserves to suffer for what he did. \newline 2.~This is payback for the harm he caused. \newline 3.~An eye for an eye --- he should feel the same pain. \\
    \addlinespace
    Suffering      & 1.~I want him to suffer in prison. \newline 2.~Prison should be painful and miserable. \newline 3.~Let him rot and suffer behind bars. \\
    \addlinespace
    Exclusion      & 1.~He has no place in our society. \newline 2.~Lock him up and throw away the key. \newline 3.~He forfeited his right to be part of society. \\
    \bottomrule
  \end{tabular}
\end{table}

\subsection{BERTopic Configuration}
\label{sec:supp_bertopic}

The unsupervised topic model was BERTopic \citep{grootendorst2022bertopic}. The pipeline used BGE embeddings as input, UMAP \citep{mcinnes2018umap} for dimensionality reduction (\texttt{n\_components} = 5, \texttt{n\_neighbors} = 15, \texttt{min\_dist} = 0.0, cosine metric), and HDBSCAN \citep{mcinnes2017hdbscan} for density-based clustering (\texttt{min\_cluster\_size} = 25, \texttt{min\_samples} = 5). Topic representation used class-based TF-IDF on the cluster-conditional vocabulary. Random seed was fixed at 42 throughout for reproducibility.

The procedure recovered seven coherent topic clusters and an outlier set of 131 responses (26.4\%) that the algorithm could not confidently assign---a typical outlier rate for short opinion text in BERTopic. The seven topics, with their automatically generated names and counts, are reported in Table~\ref{tab:supp_bertopic}. Topic numbers are arbitrary identifiers; the underlying themes are: violent details of the vignette (Topic 0), sentencing length and severity (Topic 1), Darryl as a person and his place in society (Topic 2), removal from society and crime prevention (Topic 3), rehabilitation through time and change (Topic 4), the restraining-order details specific to vignette 2 (Topic 5), and aspirations of help and hope for the future (Topic 6).

\begin{table}[htbp]
  \centering
  \caption{BERTopic clusters with automatically generated names and counts. Topic $-1$ is the outlier set ($n = 131$, 26.4\%) of responses the algorithm could not confidently assign.}
  \label{tab:supp_bertopic}
  \small
  \begin{tabular}{lrr l}
    \toprule
    Topic & $n$ & \% & Top representative tokens \\
    \midrule
    Topic 0  & 62 & 12.5 & frankie, darryl, death, frankies \\
    Topic 1  & 66 & 13.3 & sentence, years, violence, darryl \\
    Topic 2  & 67 & 13.5 & darryl, life, society, person \\
    Topic 3  & 69 & 13.9 & life, streets, society, crime \\
    Topic 4  & 26 & 5.2  & time, chance, change, rehabilitated \\
    Topic 5  & 35 & 7.1  & death, restraining order, restraining, order \\
    Topic 6  & 40 & 8.1  & life, think, help, hope \\
    Outliers & 131 & 26.4 & --- \\
    \bottomrule
  \end{tabular}
\end{table}

\subsection{Auxiliary Lexical Features (VADER, Empath)}

Two off-the-shelf lexicons were applied to each response: VADER \citep{hutto2014vader} for compound sentiment ($-1$ to $+1$ scale) and Empath \citep{fast2016empath} for category-probability scores across 26 retained content categories (e.g., violence, crime, death, suffering, healing, government, family). Both lexicons were applied to the concatenated response text. VADER scores enter the hierarchical regressions as a control feature; Empath correlations with the psychological measures are reported separately in Section~\ref{sec:supp_empath}.

\section{Study 2 Validation and Sensitivity Analyses}
\label{sec:supp_nlp_validation}

\subsection{Benchmark Validation}

To assess classifier accuracy independent of the survey responses, we constructed a 40-sentence ground-truth set: ten paradigmatic sentences for each of four core themes (deterrence, rehabilitation, retribution, revenge). Sentences were drafted to express each theme unambiguously and were not used in the prototype-similarity analyses. Each classifier was scored against the ground-truth labels on top-1 accuracy (the modal predicted label matches the ground-truth label), top-2 accuracy, and top-3 accuracy. Mean rank of the true label is reported for completeness.

\begin{table}[htbp]
  \centering
  \caption{Classifier accuracy against the 40-sentence ground-truth set.}
  \label{tab:supp_benchmark}
  \small
  \begin{tabular}{lrrrr}
    \toprule
    Classifier & Top-1 & Top-2 & Top-3 & Mean rank of true label \\
    \midrule
    Claude Haiku 4.5         & $1.000$ & $1.000$ & $1.000$ & $1.0$ \\
    DeBERTa-v3-large         & $0.875$ & $0.950$ & $0.975$ & $1.2$ \\
    BART-large-MNLI (legacy) & $0.775$ & $0.925$ & $0.950$ & $1.4$ \\
    \bottomrule
  \end{tabular}
\end{table}

The legacy BART-MNLI classifier (the model originally specified in the pre-registration) was retired from the main pipeline after this validation showed DeBERTa-v3-large achieves substantially higher accuracy on the same ground-truth set. Claude Haiku~4.5 was added as a third method and scored 100\% on the benchmark.

\subsection{Cross-Method Convergence}

Five-method pairwise convergence on the broad prosocial/dark classification is reported in Table~\ref{tab:supp_convergence}. Raw agreement is high across pairs (range $.65$ to $.95$). Cohen's $\kappa$ is modest because of the heavy class imbalance: when one class accounts for most of the data, even highly accurate classifiers can show $\kappa$ values close to zero. We report raw agreement as the headline convergence statistic in the main text, with $\kappa$ here for completeness.

\begin{table}[htbp]
  \centering
  \caption{Pairwise cross-method convergence on the broad prosocial/dark classification. Raw agreement and Cohen's $\kappa$. Sample sizes vary because the dictionary classifier left 75 responses unclassified.}
  \label{tab:supp_convergence}
  \small
  \begin{tabular}{ll r r r}
    \toprule
    Method 1 & Method 2 & $n$ & Raw agreement & $\kappa$ \\
    \midrule
    Dictionary & DeBERTa & 421 & $.72$ & $.03$ \\
    Dictionary & Claude  & 421 & $.95$ & $.12$ \\
    Dictionary & BGE     & 421 & $.73$ & $.06$ \\
    Dictionary & voyage  & 421 & $.88$ & $.09$ \\
    DeBERTa    & Claude  & 496 & $.75$ & $.16$ \\
    DeBERTa    & BGE     & 496 & $.65$ & $.13$ \\
    DeBERTa    & voyage  & 496 & $.67$ & $.02$ \\
    Claude     & BGE     & 496 & $.72$ & $.09$ \\
    Claude     & voyage  & 496 & $.87$ & $.16$ \\
    BGE        & voyage  & 496 & $.76$ & $.31$ \\
    \bottomrule
  \end{tabular}
\end{table}

\subsection{Prototype-Set Sensitivity}

The central rehabilitation-suppression finding (BGE rehabilitation similarity $\times$ punitiveness $r = -.37$) was re-tested across all 15 combinations of three embedding models (BGE, voyage, mpnet) and five prototype sets (Original, Formal, Colloquial, LayFormal, AntiTaut). Table~\ref{tab:supp_protosens} reports the rehabilitation correlation for each combination. The effect is robust: every combination yields a substantial negative correlation between rehabilitation similarity and punitiveness (range $r = -.21$ to $r = -.37$). The same pattern holds across all three embedding models even when the AntiTaut set is used, which was specifically constructed to minimize lexical overlap with the punitiveness scale items.

\begin{table}[htbp]
  \centering
  \caption{Pearson correlations between rehabilitation similarity and punitiveness across 15 combinations of embedding model and prototype set.}
  \label{tab:supp_protosens}
  \small
  \begin{tabular}{lrrrrr}
    \toprule
    Embedding model & Original & Formal & Colloquial & LayFormal & AntiTaut \\
    \midrule
    BGE-large-en-v1.5    & $-.37$ & $-.30$ & $-.34$ & $-.32$ & $-.27$ \\
    voyage-3-large       & $-.31$ & $-.28$ & $-.30$ & $-.27$ & $-.21$ \\
    Sentence-MPNet-base  & $-.28$ & $-.24$ & $-.27$ & $-.24$ & $-.22$ \\
    \bottomrule
  \end{tabular}
\end{table}

\subsection{Vignette Stability of the Facade Findings}

The central facade matrix and rehabilitation-suppression result were re-computed within each of the three Study~1 vignettes. The rehabilitation--punitiveness correlation was $r = -.36$ in Vignette~1, $r = -.32$ in Vignette~2, and $r = -.45$ in Vignette~3 (all $p < .001$). A mixed-effects model with vignette as a random intercept yielded a near-identical fixed-effect estimate to the pooled model ($\beta = -4.45$ for the regression with all four text features; $t = -8.42$, $p < .001$) and an ICC of approximately $.001$, confirming that the effect operates at the participant level and is not a between-vignette artifact.
\section{Study 2 Auxiliary Analyses}
\label{sec:supp_nlp_aux}

\subsection{Empath Lexicon Correlations With Psychological Measures}
\label{sec:supp_empath}

The Empath lexicon \citep{fast2016empath} provides probability scores for category membership across a large set of content domains. We computed Empath scores on each response and correlated 26 retained content categories with five psychological measures: punitiveness, hostile aggression, crime concerns, RWA, and racial resentment. Table~\ref{tab:supp_empath} reports correlations sorted by absolute punitiveness correlation. The pattern is consistent with the cultural-default reading: the strongest correlates of punitiveness are content categories that signal emotional engagement with violence and harm (negative\_emotion, death, kill, suffering) rather than categories that name dark motives directly. Categories tied to prosocial framing (achievement, healing, help) show negative correlations with punitiveness and hostile aggression.

\begin{table}[htbp]
  \centering
  \caption{Pearson correlations between Empath category scores and selected psychological measures. Sorted by $|r|$ with punitiveness. Twelve of 26 categories show $|r| > .10$ with punitiveness; FDR-corrected at $q = .05$ across $26 \times 5 = 130$ tests.}
  \label{tab:supp_empath}
  \footnotesize
  \begin{tabular}{lrrrrr}
    \toprule
    Empath category   & Punit. & Hostile Aggr & Crime Conc. & RWA & Racial Resent. \\
    \midrule
    negative\_emotion & $.25$  & $.13$  & $.03$  & $.03$ & $.13$ \\
    death             & $.18$  & $.10$  & $.06$  & $.09$ & $.13$ \\
    kill              & $.16$  & $.12$  & $.07$  & $.09$ & $.13$ \\
    children          & $.16$  & $.09$  & $.02$  & $.04$ & $.09$ \\
    suffering         & $.13$  & $.09$  & $.03$  & $.04$ & $.08$ \\
    achievement       & $-.11$ & $-.10$ & $-.05$ & $-.10$ & $-.08$ \\
    crime             & $.10$  & $.10$  & $.08$  & $.14$ & $.15$ \\
    violence          & $.10$  & $.07$  & $-.01$ & $.04$ & $.08$ \\
    fear              & $.08$  & $.04$  & $.03$  & $.06$ & $.04$ \\
    pain              & $.08$  & $.11$  & $.01$  & $.01$ & $.07$ \\
    positive\_emotion & $.07$  & $-.06$ & $-.03$ & $-.10$ & $.03$ \\
    fight             & $.07$  & $-.02$ & $.04$  & $.06$ & $.06$ \\
    healing           & $-.06$ & $-.13$ & $.02$  & $-.08$ & $-.10$ \\
    \multicolumn{6}{c}{\textit{(13 additional categories with $|r| < .10$ across all five measures)}} \\
    \bottomrule
  \end{tabular}

  \vspace{0.3em}
  \footnotesize
  \textit{Note.} Full 26-category correlation table available in the project repository (\texttt{empath\_correlations\_full.csv}).
\end{table}

\subsection{BERTopic Distribution by Hostile-Aggression Tertile and Political Group}
\label{sec:supp_bertopic_by_group}

The seven BERTopic clusters were cross-tabulated against hostile-aggression tertile (low / mid / high) and against political group (liberal / moderate / conservative) to test the cultural-default reading directly. The omnibus $\chi^2$ tests were significant at conventional thresholds, $\chi^2(12) = 21.38$, $p = .045$ for tertile and $\chi^2(12) = 23.47$, $p = .024$ for political group. Cell-level differences were small, with no topic appearing in fewer than 3\% or more than 31\% of responses in any subgroup, and the rank ordering of topic frequencies was preserved across all three political groups.

The single largest cell-level difference emerged for Topic~4 (rehabilitation theme): 8.0\% of low-hostile responses fell into this topic versus 2.4\% of high-hostile responses, a 3.3-fold difference. This complements the rehabilitation-suppression finding from the prototype-similarity analyses: the unsupervised topic model independently recovers the rehabilitation theme as the one whose representation differs most across the hostility distribution.

\subsection{Political Moderation Tests}
\label{sec:supp_political_mod}

To test whether the central facade-matrix patterns differ by political orientation, we ran moderated multiple regressions of the form $Y = \beta_0 + \beta_1 X + \beta_2 \text{Polit} + \beta_3 (X \times \text{Polit}) + \epsilon$, where $X$ is the relevant text feature and $Y$ is the relevant psychological measure. Eight interaction tests were conducted, one per facade-matrix cell of theoretical interest (the four strongest negative cells $\times$ two political moderators, where the moderator was either the continuous 7-point political-orientation scale or a dummy contrast between liberals and conservatives).

After FDR correction across the eight tests, no interaction reached significance ($p_\textit{FDR}$ values ranged from $.21$ to $.92$). The rehabilitation-suppression effect, the Claude prosocial--dark gap effect, and the BGE prosocial--dark gap effects all hold across the political spectrum. The cultural default operates similarly for liberals and conservatives.

\section{Effect-Size Magnitude Comparisons}
\label{sec:supp_effect_sizes}

For readers who want a single comparison-of-magnitudes summary, Table~\ref{tab:supp_effect_summary} reproduces the main findings of both studies in a common metric: Pearson $r$ where applicable, with $\Delta R^2$ for the hierarchical regressions. Effects of practical importance ($|r| > .10$ or $\Delta R^2 > .02$) are flagged.

\begin{table}[htbp]
  \centering
  \caption{Summary of central effect sizes from both studies.}
  \label{tab:supp_effect_summary}
  \footnotesize
  \begin{tabular}{lll}
    \toprule
    Effect & Magnitude & Source \\
    \midrule
    \textit{Study 1: Confirmatory} \\
    Punitiveness $\times$ hostile aggression cluster   & $r = .59$ & Main Table~\ref{tab:H1_correlations} \\
    Punitiveness $\times$ emotions cluster             & $r = .53$ & Main Text \\
    Punitiveness $\times$ personality \& ideology cluster & $r = .47$ & Main Text \\
    Punitiveness $\times$ crime concerns cluster       & $r = .32$ & Main Table~\ref{tab:H1_correlations} \\
    Hostile aggression vs.\ crime concerns difference  & $\Delta r = .26$, $Z = 6.26$ & Main Table~\ref{tab:steiger} \\
    \addlinespace
    \textit{Study 1: Sensitivity} \\
    Tautology sensitivity (Original 6-construct cluster)  & $\Delta r = .26$ & Section~\ref{sec:supp_bootstrap} \\
    Tautology sensitivity (Strictest 2-construct cluster) & $\Delta r = .22$ & Section~\ref{sec:supp_bootstrap} \\
    Sentence anchor: pooled departure from 25-year midpoint & $d = 0.81$ & Section~\ref{sec:supp_anchor} \\
    \addlinespace
    \textit{Study 2: Facade Matrix} \\
    BGE rehabilitation $\times$ punitiveness               & $r = -.37$ & Main Table~\ref{tab:nlp_facade} \\
    Claude prosocial--dark $\times$ punitiveness           & $r = -.28$ & Main Table~\ref{tab:nlp_facade} \\
    Claude prosocial--dark $\times$ hostile aggression     & $r = -.25$ & Main Table~\ref{tab:nlp_facade} \\
    BGE rehabilitation $\times$ RWA                        & $r = -.20$ & Main Table~\ref{tab:nlp_facade} \\
    BGE rehabilitation $\times$ racial resentment          & $r = -.16$ & Main Table~\ref{tab:nlp_facade} \\
    BGE rehabilitation $\times$ hostile aggression         & $r = -.15$ & Main Table~\ref{tab:nlp_facade} \\
    \addlinespace
    \textit{Study 2: Multiverse Verdicts} \\
    Cultural-default verdicts                              & 4 of 6 & Main Table~\ref{tab:nlp_multiverse} \\
    Individual-facade verdicts                             & 2 of 6 & Main Table~\ref{tab:nlp_multiverse} \\
    Sincerity verdicts                                     & 0 of 6 & Main Table~\ref{tab:nlp_multiverse} \\
    \addlinespace
    \textit{Study 2: Hierarchical Regression Increments} \\
    $\Delta R^2$, sentencing decision (text features over Study~1 scales)   & $.134$ & Main Table~\ref{tab:nlp_regression} \\
    $\Delta R^2$, punitiveness aggregate (text features over Study~1 scales) & $.079$ & Main Table~\ref{tab:nlp_regression} \\
    \bottomrule
  \end{tabular}
\end{table}

\section{Reproducibility}

All analysis code, intermediate features, and final analytical datasets are publicly available at \url{https://github.com/dgk-law-and-cognition-lab/Punishment-Punitiveness} and at the project page \url{https://dgk-law-and-cognition-lab.github.io/Punishment-Punitiveness/}. The Python NLP pipeline runs on a fixed random seed (42) at every stochastic step. The R analysis scripts are organized into three Rmd notebooks: \texttt{01\_main\_analysis.Rmd} (confirmatory tests), \texttt{02\_nlp\_integration.Rmd} (linking text features to psychological measures), and \texttt{03\_supplementary\_analyses.Rmd} (the analyses reported in this appendix). Package versions and random seeds are recorded in the session-info output of each notebook.
